\section{Experimental Setup}
\label{sec:experiments}

\subsection{Datasets}
The study is designed around four complementary climate/weather sources: Central England Temperature (CET), NOAA/NCEI GHCN-Daily station observations, Jena weather-station measurements, and ERA5 reanalysis. CET provides continuity with prior climate KAN work; GHCN-Daily evaluates station-level geographic generalization; Jena provides high-frequency multivariate observations; and ERA5 supports reanalysis-based multivariate and distribution-shift experiments. Exact station choices, variables, date ranges, units, missing-data rules, licenses, and download procedures must be frozen before headline experiments. \todo{Populate Table~\ref{tab:datasets} only after dataset metadata are finalized.}

\begin{table}[t]
\centering
\caption{Dataset characteristics. All entries must be generated from finalized dataset metadata.}
\label{tab:datasets}
\begin{tabular}{lcccc}
\toprule
Dataset & Resolution & Variables & Period & Samples \\
\midrule
CET & \todo{daily} & \todo{} & \todo{} & \todo{} \\
GHCN-Daily & \todo{} & \todo{} & \todo{} & \todo{} \\
Jena & \todo{} & \todo{} & \todo{} & \todo{} \\
ERA5 & \todo{} & \todo{} & \todo{} & \todo{} \\
\bottomrule
\end{tabular}
\end{table}

\subsection{Chronological Splitting and Leakage Prevention}
All time series are split chronologically into training, validation, calibration, and final test regions. Normalization parameters are estimated from training observations only. Windowed samples are assigned according to forecast-target origin, and windows whose forecast target crosses a split boundary are excluded. Hyperparameter selection uses training and validation data; conformal correction and reliability thresholds use calibration data; the final test set is reserved exclusively for evaluation.

\subsection{Forecasting Tasks}
Experiments include multiple forecast horizons selected according to each dataset's temporal resolution. The CET development benchmark currently considers 1-, 7-, 30-, and 90-step horizons. Equivalent horizon sets for other datasets will be specified before final experiments so that comparisons are meaningful within each temporal resolution.

\subsection{Baselines}
The benchmark includes persistence and classical machine-learning models, recurrent networks, temporal convolution, Transformer and state-space sequence models, a standard KAN baseline, a Tem$^2$-KAN reference implementation, and the proposed \method. Optional models are reported as failed/unavailable rather than silently removed if dependencies or execution fail. Hyperparameter tuning effort should be comparable across major learned baselines.

\subsection{Evaluation Metrics}
Deterministic forecasting is evaluated using MAE and RMSE in physical target units. Probabilistic forecasting uses empirical interval coverage, mean prediction-interval width, and interval/calibration scores where supported. Reliability is evaluated using reliability--error correlation, large-error detection AUROC/AUPRC, risk--coverage curves, AURC, retained-sample error, and abstention rate. Efficiency analysis reports parameter count, training time, inference latency, and memory where measurable.

\subsection{Statistical Analysis}
Headline neural comparisons use at least five independent random seeds. Paired comparisons are computed from predictions at identical forecast origins and timestamps. Because adjacent and multi-horizon windows overlap, primary confidence intervals for MAE and RMSE differences use a circular moving-block bootstrap over forecast origins, keeping all horizon elements from each origin together. Block length is fixed before final comparison and is at least the forecast-horizon width. Origin-level Wilcoxon tests and paired standardized mean differences are reported only as sensitivity analyses because they do not model serial dependence. Holm correction is applied to each pre-specified primary comparison family; exploratory families may additionally report Benjamini--Hochberg correction. Statistical evidence is interpreted together with effect magnitude and practical relevance rather than as a binary criterion.

\subsection{Robustness and Shift Protocols}
Robustness experiments include sensor noise, missing inputs, temporal regime shift, and extreme-event subsets. Adaptive/rolling conformal experiments evaluate coverage through time rather than only as an aggregate. Any synthetic corruption level or shift definition must be fixed before examining final model rankings.
